%% The normal-projected semi-Lagrangian (nSL) level-set method: derivation,
%% implementation, and first verification results (including the measured
%% negative results that bound its present form).
%% Build:  latexmk -pdf normalProjectedSemiLagrangian.tex
\documentclass[preprint,3p,11pt]{elsarticle}

\usepackage[T1]{fontenc}
\usepackage{amsmath,amssymb}
\usepackage{bm}
\usepackage{booktabs}
\usepackage{url}
\usepackage[hidelinks]{hyperref}

% --- notation ----------------------------------------------------------------
\newcommand{\psilev}{\psi}
\newcommand{\xc}{\mathbf{x}_c}
\newcommand{\xd}{\mathbf{x}_d}
\newcommand{\uvec}{\mathbf{u}}
\newcommand{\un}{\mathbf{u}_n}
\newcommand{\nv}{\mathbf{n}}
\newcommand{\gv}{\mathbf{g}}
\newcommand{\Hm}{\mathbf{H}}
\newcommand{\gradpsi}{\lvert\nabla\psilev\rvert}

\journal{a computational physics journal (method-line note)}

\begin{document}
\begin{frontmatter}

\title{A normal-projected semi-Lagrangian level-set update with per-step
geometric renormalization: derivation, implementation, and measured limits}

\author[tud,fe]{Tomislav Mari\'c}
\address[tud]{Mathematical Modeling and Analysis, Technische Universit\"at
Darmstadt, Germany}
\address[fe]{Focused Energy, Darmstadt, Germany}

\begin{abstract}
The baseline semi-Lagrangian level-set transport
$\psilev^{n+1}(\xc)=\psilev^{n}(\xd)$ resamples the transported field through a
per-cell quadratic reconstruction at every step; on rigid-translation tests this
reconstruct-and-evaluate operator amplifies a grid-scale ($\lambda\approx 2h$)
corrugation of the zero set at a rate that is fixed per unit interface
displacement. The normal-projected semi-Lagrangian update (nSL) examined here
replaces the value resampling by a purely geometric update: the departure foot
is traced along the normal-projected velocity $\un=(\uvec\cdot\nv)\nv$ -- exact
for the continuous problem, since the tangential velocity carries no information
about $\psilev$ -- and the written value is composed of the cell's geometric
offset from its own local model's zero set plus the normal displacement of the
trace, renormalizing $\gradpsi\to1$ in the interface band at every step without
any global redistancing operator. We give the complete discrete formulation:
the departure-centred second-order trace built from the current and previous
projected velocities, the analytic assembly of $\nabla\un$ from the fit, the
stable quadratic-root offset conversion with its acceptance guard, the
stabilized foot-point algorithm as an alternative curvature-aware offset
engine, and a strain-rescaling alternative that integrates the exact
gradient-magnitude transport equation. The implementation is a
runtime-selectable advection scheme in OpenFOAM with every control a validated
dictionary entry. The verification record is reported in full, including its
negative results: the trace itself is clean and preserves an exact free stream
to machine precision, but every per-step \emph{write-back} of a fit-derived
offset diverges at an engine-independent rate -- the quadratic root, a
first-order conversion, and the stabilized foot-point produce the same band
error growth -- and the grid-scale corrugation amplification is unchanged
relative to the baseline, because the noise re-enters through the fitted
normals rather than through the value resampling. The measured gates that
bound each design choice are stated so that subsequent variants can be judged
against them.
\end{abstract}

\begin{keyword}
level set \sep semi-Lagrangian \sep signed distance \sep renormalization
\sep foot point \sep OpenFOAM
\end{keyword}
\end{frontmatter}

% =============================================================================
\section{Motivation}
\label{sec:motivation}

Two measured properties of the baseline quadratic semi-Lagrangian transport
motivate this method line. First, on an essentially exact rigid translation
(coupled solver, surface tension coefficient $\sigma=0$, velocity a uniform
free stream to $\sim\!3\times10^{-15}$\,m/s), the $\psilev=0$ contour develops
a high-wavenumber corrugation ($\lambda\approx2h$, azimuthal modes $m>4$) that
grows exponentially at a rate fixed per unit interface displacement
($\approx0.19$--$0.28$ per cell of travel), essentially independent of the
Courant number over $\mathrm{Co}\in[0.007,0.5]$. Second, in the two-phase
coupling the band gradient magnitude $\gradpsi$ orders every observed
stationary/translating/oscillating-droplet outcome, and its collapse
\emph{precedes} the parasitic-current growth. The nSL design targets both: it
removes the per-step polynomial resampling of $\psilev$ from the value path,
and it renormalizes the band toward a signed distance at every step by a
purely local, zero-set-preserving construction -- explicitly avoiding global
redistancing, which was measured (in this framework) to rebuild distance to
advection artefacts.

% =============================================================================
\section{Mathematical formulation}
\label{sec:method}

\subsection{Continuous foundation: the tangential velocity is redundant}
\label{sec:exact}

Decompose the velocity at any point into components normal and tangential to
the local level set of $\psilev$,
\begin{equation}
  \uvec = v_n\,\nv + \uvec_t,
  \qquad v_n=\uvec\cdot\nv,
  \qquad \nv=\frac{\nabla\psilev}{\gradpsi}.
\end{equation}
Because $\nabla\psilev$ is normal to its own level sets,
$\uvec_t\cdot\nabla\psilev=0$ \emph{identically}, so
\begin{equation}
  \uvec\cdot\nabla\psilev
  = v_n\,\gradpsi
  = \un\cdot\nabla\psilev,
  \qquad \un := v_n\,\nv,
\end{equation}
and the level-set equation is equivalent to
\begin{equation}
  \frac{\partial\psilev}{\partial t} + \un\cdot\nabla\psilev = 0 .
  \label{eq:nsl-pde}
\end{equation}
$\psilev$ is therefore exactly constant along trajectories of $\un$, for any
field, drifted or not: tracing backward along $\un$ solves the same continuous
transport problem as tracing along $\uvec$, and the tangential velocity is
discarded without loss. Two consequences fix the design: pure tangential
motion is an exact no-op ($\un=\mathbf 0\Rightarrow\xd=\xc$), and the update
never needs the tangential component of the flow.

\subsection{The discrete trace}
\label{sec:trace}

Per cell $c$ with centre $\xc$, the quadratic weighted least-squares fit of
$\psilev^{n}$ over the cell's stencil (identical to the baseline
reconstruction) supplies the value $\psilev_c$, the gradient
$\gv=\nabla R_c(\xc)$ and the constant Hessian $\Hm=\nabla\nabla R_c$ of the
local model $R_c$. The unit normal is $\nv=\gv/\lvert\gv\rvert$, and the
normal-projected velocity at the current level is
\begin{equation}
  \un^{\,n} = \bigl(\uvec^{\,n}\cdot\nv\bigr)\,\nv .
\end{equation}

The coupled solver advances the interface \emph{before} its momentum solve, so
only $\uvec^{\,n}$ and the stored projection $\un^{\,n-1}$ are available. The
departure foot is therefore the \emph{departure-centred} second-order backward
trace -- the only centring whose normals exist at their own time level:
\begin{equation}
  \boxed{\;
  \xd=\xc
  -\Delta t\Bigl[\un^{\,n}
      +\tfrac12\,r\,\bigl(\un^{\,n}-\un^{\,n-1}\bigr)\Bigr]
  +\tfrac12\,\Delta t^{2}\,\bigl(\un^{\,n}\!\cdot\nabla\bigr)\un^{\,n},
  \qquad r=\frac{\Delta t}{\Delta t_0} \; }
  \label{eq:foot}
\end{equation}
which at uniform step ($r=1$) is
$\xd=\xc-\tfrac{\Delta t}{2}\,(3\un^{\,n}-\un^{\,n-1})
+\tfrac12\Delta t^2(\un^{\,n}\cdot\nabla)\un^{\,n}$: the linear displacement
uses the time-centred midpoint velocity, obtained by Adams--Bashforth
extrapolation of the two available levels. $\un^{\,n-1}$ is the projected
velocity \emph{stored at the end of the previous step, projected with that
step's normals}; re-projecting the old velocity with the new normals mixes
time levels of the geometry and is not equivalent.

The velocity gradient of the projected field, needed by the convective term,
is assembled \emph{analytically} from the fit and the solver-supplied
$\nabla\uvec$ -- never by differencing the projected field. With the
convention $(\nabla\uvec)_{ij}=\partial_i u_j$ (OpenFOAM's
\texttt{fvc::grad}; the transpose of the convention
$\mathrm d\uvec=(\nabla\uvec)\,\mathrm d\mathbf x$):
\begin{align}
  \partial_i n_j &= \frac{H_{ij}-(\Hm\nv)_i\,n_j}{\lvert\gv\rvert},
  \label{eq:gradn}\\
  \partial_i v_n &= (\nabla\uvec\cdot\nv)_i+(\nabla\nv\cdot\uvec)_i,\\
  \partial_i (u_n)_j &= (\nabla v_n)_i\,n_j + v_n\,\partial_i n_j,
\end{align}
and $\bigl(\un\cdot\nabla\bigr)\un$ is the contraction
$(u_n)_i\,\partial_i(u_n)_j$. Equation~\eqref{eq:gradn} follows from
$\nabla R_c=\gv+\Hm\mathbf d$ and the quotient rule; for a signed-distance
field it reduces to the classical $\nabla\nv=(\mathbf I-\nv\nv^{\mathsf T})
\Hm/\lvert\gv\rvert$ written with the projector on the appropriate side of the
symmetric $\Hm$.

The trace obeys the baseline stencil restriction
$\lvert\xd-\xc\rvert\le$ stencil radius (warn-only guard); since
$\lvert\un\rvert\le\lvert\uvec\rvert$ pointwise, the nSL trace is never longer
than the baseline trace.

\subsection{The geometric update and the offset conversion}
\label{sec:update}

The update writes
\begin{equation}
  \boxed{\;\psilev^{n+1}(\xc)=d_c+\delta,\qquad
  \delta=(\xd-\xc)\cdot\nv\;}
  \label{eq:update}
\end{equation}
where $d_c$ is the signed geometric offset from $\xc$ to the zero set of the
cell's \emph{own} local model, and $\delta$ is the normal projection of the
trace displacement (the arc-length alternative differs at
$O(\Delta t^{3})$ and is deferred). No value of $\psilev$ is interpolated at
$\xd$: the reconstruction is consulted for directions, speeds and geometry
only. Sign conventions: $\nv$ points toward increasing $\psilev$; $d_c$
carries the sign of $\psilev_c$; a flow moving the interface in the $+\nv$
direction has $v_n>0$ and $\delta<0$. Limiting cases: $v_n=0$ gives
renormalization only; on a clean signed distance with uniform $v_n$,
$d_c+\delta=\psilev_c-v_n\Delta t$ exactly.

\paragraph{Offset engine (b): the stable quadratic root.}
Restricting the model to the normal ray, $R(s)=\psilev_c+s\lvert\gv\rvert
+\tfrac12 s^2 h_{nn}$ with $h_{nn}=\nv^{\mathsf T}\Hm\,\nv$, the signed
distance to the model's zero level is, in cancellation-free form,
\begin{equation}
  d_c=\frac{2\,\psilev_c}{\;\lvert\gv\rvert+\sqrt{D}\;},
  \qquad D=\lvert\gv\rvert^{2}-2\,\psilev_c\,h_{nn},
  \label{eq:root}
\end{equation}
taking the root that reduces to $\psilev_c/\lvert\gv\rvert$ as $h_{nn}\to0$.
It is exactly zero-preserving ($\psilev_c=0\Rightarrow d_c=0$ for any
$\gv,\Hm$): the renormalization never moves the front. The root is accepted
while $D\ge\beta\lvert\gv\rvert^{2}$ (default $\beta=\tfrac14$) -- comfortably
before the model's turning point, where the conversion becomes
ill-conditioned -- and replaced by the first-order $\psilev_c/\lvert\gv\rvert$
otherwise. The identical conversion serves the curvature path
(\texttt{offsetCorrection quadraticRoot}), so transport and curvature never
disagree about a cell's offset.

\paragraph{Offset engine (a): the stabilized foot-point.}
The curvature-aware alternative computes both legs of \eqref{eq:update} as
foot-point distances on the local model: $d_c$ as the signed distance of
$\xc$ to $\{R_c=0\}$, and the displacement as the signed distance of $\xd$ to
the cell's own iso-surface $\{R_c=\psilev_c\}$, replacing the flat projection
$\delta$. The foot point of a query $\mathbf p$ on $\{G=0\}$,
$G(\mathbf d)=R_c(\xc+\mathbf d)-\text{level}$, is found by alternating two
projections -- the tangent-plane foot point at the current surface iterate and
a Newton \texttt{surfacepoint} pull back to the surface -- accelerated by a
parabola correction: with $\mathbf f_1$ the tangential step and $\mathbf f_2$
the measured pull-back of one cycle, the curve
$\mathbf x(\alpha)=\mathbf p_i+\alpha\mathbf f_1+\alpha^2\mathbf f_2$ is a
local quadratic model of the surface path, the minimizer of
$\lvert\mathbf p-\mathbf x(\alpha)\rvert^2$ satisfies a cubic whose one Newton
step from $\alpha=1$ is
\begin{equation}
  \alpha=1-\frac{a_0+a_1+a_2+a_3}{a_1+2a_2+3a_3},
  \quad
  \begin{aligned}
  a_0&=(\mathbf p-\mathbf p_i)\cdot\mathbf f_1, &
  a_1&=2\mathbf f_2\cdot(\mathbf p-\mathbf p_i)-\lvert\mathbf f_1\rvert^2,\\
  a_2&=-3\,\mathbf f_1\cdot\mathbf f_2, &
  a_3&=-2\lvert\mathbf f_2\rvert^2,
  \end{aligned}
\end{equation}
accepted inside a guard interval $0<\alpha<\alpha_{\max}$. The signed distance
is $(\mathbf p-\mathbf x^{\Sigma})\cdot\nv^{\Sigma}$ with the sign along the
model gradient at the foot. All evaluations are polynomial arithmetic in the
fit coefficients; the foot must remain inside the stencil trust region. The
complete algorithm, its failure modes and tolerances are documented in the
companion note \texttt{stable-foot-point-3d.md}.

\subsection{The update ladder: where each branch applies, and why}
\label{sec:ladder}

Each gate below was forced by a measured failure; the numbers are in
Section~\ref{sec:results}.
\begin{enumerate}
\item \textbf{Band} ($\lvert\psilev_c\rvert/\lvert\gv\rvert\le
  k_{\mathrm{band}}\,\times$ stencil radius, default $k_{\mathrm{band}}=1$):
  the renormalizing update \eqref{eq:update}. The root \eqref{eq:root} is
  meaningful only where the local fit \emph{samples} the zero set it converts
  to: beyond one stencil radius its $\psilev_c h_{nn}$ term amplifies
  fitted-Hessian noise linearly in the offset, and cells whose stencils reach
  a level-set skeleton (a circle's cone tip) seed a closed-loop runaway.
\item \textbf{Far field}: raw first-order transport along the ray,
  $\psilev^{n+1}=\psilev_c+\delta\,\lvert\gv\rvert$ -- no division by the
  fitted gradient anywhere, so fit error enters only through the per-step
  increment, and a clamped-plateau cell ($\lvert\gv\rvert\approx0$) freezes
  automatically. Any far-field \emph{renormalization} divides by the fitted
  $\lvert\gv\rvert$ and amplifies its error linearly in the offset.
\item \textbf{Degenerate cells} (constant-fit fallback, or
  $\lvert\gv\rvert<$ \texttt{minGradPsi}): frozen for the step,
  $\psilev^{n+1}=\psilev_c$; they re-activate when the geometry regularizes.
\end{enumerate}

\subsection{Strain renormalization: the $O(\Delta t)$-gain alternative}
\label{sec:strain}

The exact transport of the gradient magnitude along trajectories,
$\mathrm D\gradpsi/\mathrm Dt=-\gradpsi\,\varepsilon_{nn}$ with
$\varepsilon_{nn}=\nv\cdot\mathbf D\cdot\nv$ and
$\mathbf D=\tfrac12(\nabla\uvec+\nabla\uvec^{\mathsf T})$, implies that
band-scale distances stretch by $(1+\varepsilon_{nn}\Delta t)$ per step. The
\texttt{strain} mode therefore writes, in the band,
\begin{equation}
  \psilev^{n+1}(\xc)
  =\bigl(\psilev_c+\delta\,\lvert\gv\rvert\bigr)
   \bigl(1+\varepsilon_{nn}\,\Delta t\bigr),
  \label{eq:strain}
\end{equation}
with raw transport outside. Its decisive structural property is the noise
gain: $\varepsilon_{nn}$ comes from the solver's smooth $\nabla\uvec$, the fit
enters only through $\nv$, and the whole correction is multiplied by
$\Delta t$ -- unlike the geometric write-back \eqref{eq:update}, whose
fit-derived offset carries no factor of $\Delta t$. The band gate on
\eqref{eq:strain} is essential: the ODE describes the stretching of a
band-scale distance whose ray foot experiences $\varepsilon_{nn}$; applied to
the large far-field $\psilev$ with the \emph{local} $\varepsilon_{nn}$, the
factor compounds $\exp(\int\varepsilon\,\mathrm dt)$ over the run.

% =============================================================================
\section{Implementation}
\label{sec:implementation}

The scheme is the runtime-selectable class \texttt{normalProjectedScheme}
(\texttt{TypeName "normalProjected"}) in the \texttt{slScheme} hierarchy of
\texttt{libleiaLevelSet}, selected by
\texttt{levelSet\{semiLagrangian\{scheme normalProjected;\}\}} in
\texttt{fvSolution}; the default \texttt{pointValue} keeps every existing case
dictionary valid. The scheme is handed the reconstruction and corrector owned
by \texttt{slAdvection}; the corrector is deliberately unused (there is no
reconstruction evaluation at the foot to correct, and its stencil-bounds clip
would be wrong for a renormalized value), and the foot-radius and non-finite
guards are applied inline.

\paragraph{Reconstruction interface.} Three optional capabilities were added
to the abstract \texttt{slReconstruction}, following the library's
\texttt{NotImplemented}-default pattern (no \texttt{dynamic\_cast} anywhere):
\texttt{fitDerivatives(c, g, H)} returning the per-cell fit order actually
available (2 full quadratic, 1 linear fallback with $\Hm=\mathbf 0$, 0
constant fallback -- the scheme's degradation ladder for free);
\texttt{signedOffset(c, fallback)} implementing \eqref{eq:root} with its
guard; and \texttt{footPointDistance(c, p, level, ok)} implementing the
stabilized foot-point on the local model. All are overridden by the
production \texttt{uncachedQuadraticWeightedLeastSquares} reconstruction and
read exclusively from the already-fitted coefficient buffer.

\paragraph{Dictionary controls.} Everything is a validated runtime entry in
\texttt{levelSet.semiLagrangian}; nothing is hardcoded:
\begin{center}\small
\begin{tabular}{llp{0.52\linewidth}}
\toprule
entry & default & meaning \\
\midrule
\texttt{scheme} & \texttt{pointValue} & \texttt{normalProjected} selects nSL \\
\texttt{renormalization} & \texttt{strain} & \texttt{geometric} $\vert$
  \texttt{strain} $\vert$ \texttt{none} (Sections~\ref{sec:update},
  \ref{sec:strain}) \\
\texttt{offsetEngine} & \texttt{rayRoot} & \texttt{rayRoot} $\vert$
  \texttt{footPoint}; read in \texttt{geometric} mode only \\
\texttt{bandRadii} & 1.0 & band gate in stencil radii
  (Section~\ref{sec:ladder}) \\
\texttt{minGradPsi} & 0.05 & freeze threshold on the fitted
  $\lvert\gv\rvert$ \\
\texttt{offsetBeta} & 0.25 & root acceptance threshold $\beta$
  in \eqref{eq:root} \\
\texttt{footPointTolRel} & $10^{-6}$ & foot-point tolerance, relative to the
  stencil radius \\
\texttt{footPointAlphaMax} & 20 & parabola-extrapolation guard interval \\
\texttt{footPointCycles} & 10 & cap on tangent-plane/surfacepoint cycles \\
\texttt{footPointSurfIters} & 20 & cap on \texttt{surfacepoint} Newton steps \\
\bottomrule
\end{tabular}
\end{center}
The study pipeline exposes \texttt{scheme} as the token \texttt{SL\_SCHEME}
(and the curvature conversion as \texttt{SL\_OFFSET\_CORRECTION}), so the
scheme is a first-class study axis.

\paragraph{Time levels and storage.} \texttt{advance(psi, Unew, Uold, recon,
corrector)} deliberately ignores \texttt{Unew}: the departure-centred trace
\eqref{eq:foot} needs only the time-$n$ velocity (the \texttt{Uold} argument
at every call site) and the stored projection of the previous step. That
projection lives in the registered field \texttt{uProjOld.nSL}
(\texttt{READ\_IF\_PRESENT}/\texttt{AUTO\_WRITE}), so a restarted trace uses
the same two levels as an uninterrupted run; on a cold start (and per cell on
re-activation after a freeze) the extrapolation degenerates to the identity --
the standard one-off Adams--Bashforth startup within the second-order budget.

\paragraph{Update discipline and parallelism.} The fit coefficients are
increments about the \emph{live} $\psilev[c]$, so the per-cell loop assembles
$\psilev^{n+1}$ into a scratch field and assigns it once at the end, followed
by \texttt{correctBoundaryConditions()} and the reconstruction's
\texttt{postAdvect} hook. The communication pattern is unchanged from the
baseline: one stencil halo of $\psilev$ (inside \texttt{recon.update}) plus
the standard halo of \texttt{fvc::grad(Uold)}; everything else is cell-local,
so the scheme is dimension-independent (one implementation covers 2D
one-cell-thick and 3D meshes; the fitted $\gv,\Hm$ have identically zero
components in empty directions) and embarrassingly parallel. Per-step counters
(frozen, far-field, conversion fallback, foot-point fallback, feet outside the
stencil, non-finite resets) are reduced and reported every step.

% =============================================================================
\section{Measured verification record}
\label{sec:results}

All transport measurements use the $\sigma=0$ rigid-translation rig (coupled
solver, water/air droplet geometry, $N=64$, drop radius $R=1$\,mm $=6.4h$,
$\mathbf U=(0.05,0,0)$\,m/s): with the capillary force identically zero the
velocity remains an exact free stream ($\max\lvert\mathbf U-\mathbf
U_{\mathrm{trans}}\rvert\approx3\times10^{-15}$\,m/s in every run quoted), so
every interface error is attributable to the transport scheme alone.

\paragraph{The trace is clean.} One full step of the update against the exact
translated distance field: $0.017h$ maximum error in the band, $0.027h$
everywhere. Over 1600 steps (16$h$ of travel) the front radius, phase volume,
zero-set error and band $\gradpsi$ match the \texttt{pointValue} baseline
(volume error $1.17\times10^{-2}$ vs $1.02\times10^{-2}$; band $\gradpsi$
$[0.82,1.38]$ vs $[0.84,1.37]$; zero-set $L_2$ $8.7\times10^{-5}$ vs
$7.5\times10^{-5}$\,m).

\paragraph{The geometric write-back diverges, engine-independently.} Band
error after 100 steps:
\begin{center}\small
\begin{tabular}{lc}
\toprule
in-band update & band error \\
\midrule
geometric, quadratic-root engine \eqref{eq:root} & $1.107\,h$ \\
geometric, first-order conversion $\psilev_c/\lvert\gv\rvert$ & $1.11\,h$ \\
geometric, stabilized foot-point engine & $1.113\,h$ \\
raw transport (no renormalization) & $0.012\,h$ \\
\bottomrule
\end{tabular}
\end{center}
Three offset engines of very different accuracy, one divergence rate
($\times1.7$ per 10 steps, independent of $\Delta t$): writing \emph{any}
fit-derived offset back into $\psilev$ feeds neighbour noise into the next
step's fit at order one per step
($\delta d_c\sim\psilev_c\,\delta\lvert\gv\rvert/\lvert\gv\rvert^{2}$ carries
no factor of $\Delta t$), and the offset engine's accuracy is irrelevant to
the loop gain. The foot-point engine was genuinely engaged (0--2 per-cell
fallbacks per step out of $\approx$112 band cells). Two further amplification
channels were isolated on the way and are now gated
(Section~\ref{sec:ladder}): offset-linear Hessian-noise amplification
($\sim$30$h$ one-step errors at 45$h$ offset, domain-corner cells) and
skeleton-stencil seeding (runaway from $2.6h$ off the circle's cone tip even
when the root was gated at three stencil radii).

\paragraph{The corrugation hypothesis is falsified.} The design expectation
was that removing the value resampling removes the $\lambda\approx2h$
amplifier. Measured $m>4$ corrugation at $16h$ displacement:
\texttt{pointValue} $0.209h$, nSL $0.223h$ -- the same exponential growth.
The corrugation re-enters through the \emph{fitted normals}: a grid-scale
wiggle in $\psilev$ perturbs $\gv$, hence $\nv$, hence both the trace
direction and the written increment $\delta$, with the same per-displacement
gain as the resampling it replaced. The pre-registered acceptance rule
(closed-loop amplification $\le1$ at $\lambda\approx2h$) is therefore not
met.

\paragraph{Deforming flow.} On the reversed single (shear) vortex at $T=2$
the strain-mode nSL is anti-convergent over $N=32..256$ (shape-error order
$-1.18$ at $\mathrm{Co}=0.5$, volume errors up to $O(1)$, band minimum
$\gradpsi\to0$), while the paired \texttt{pointValue} rows reproduce the
committed baseline to all printed digits. Variant (a) at $N=64$,
$\mathrm{Co}=0.5$: shape error $2.28\times10^{-2}$ against the baseline's
$4.29\times10^{-4}$. An earlier, unscoped application of \eqref{eq:strain} to
the whole domain compounded $\exp(\int\varepsilon\,\mathrm dt)$ in the far
field (band gradient errors $10^{6}$--$10^{34}$ across the ladder) and is the
measured reason for the band gate on the strain factor.

\paragraph{Status.} The value path -- the full quadratic resampling at the
foot -- remains the most accurate transport measured in this framework; every
geometric substitute tested so far is worse on deforming flows, and every
per-step renormalizing write-back diverges regardless of the offset engine.
What survives and is retained: the normal-projected trace itself (exact
free-stream preservation, baseline-equal transport on translation), the
shared second-order offset conversion for the \emph{curvature} path
(\texttt{offsetCorrection quadraticRoot}, valid where the fit samples the
zero set), the stabilized foot-point as a reusable geometric kernel, and the
strain mode as the only renormalization whose noise gain carries a factor of
$\Delta t$. The next configuration in this line applies the strain factor
\eqref{eq:strain} on top of the \emph{unchanged} baseline value transport --
one change against the verified scheme.

% =============================================================================
\section*{Code and data availability}
Implementation: \texttt{src/leiaLevelSet/semiLagrangian/normalProjectedScheme.\{H,C\}},
reconstruction extensions in
\texttt{uncachedQuadraticWeightedLeastSquaresReconstruction.\{H,C\}}; design
notes and the foot-point companion in
\texttt{docs/normal-projected-semi-lagrangian/}; study configurations
\texttt{config/nsl*.yaml}, \texttt{config/npsl*.yaml}. Repository:
\url{https://github.com/leia-openfoam/leia}.

\end{document}
